\textbf{(CAP-COLUNI 2020 -- Q15)} Um \textit{smartphone} apresenta, na parte superior de sua tela, dois indicadores. Um deles refere-se ao sinal do \textit{wi-fi}, representado por um setor circular, enquanto o outro, representado por um retângulo, indica a quantidade de bateria do aparelho. Na Figura 1 abaixo, são apresentadas as formas e as dimensões geométricas dos indicadores de sinal de \textit{wi-fi} e de quantidade de bateria, respectivamente.

\includegraphics{./figs/2020-Q15-A.png}

Na Figura 2, estão apresentadas situações nas quais o sinal de \textit{wi-fi} está muito bom, portanto o setor circular encontra-se completamente preenchido (hachurado), enquanto a bateria do celular apresenta certo desgaste, por isso o retângulo não está completamente preenchido.

\includegraphics{./figs/2020-Q15-B.png}

Considere a situação na qual o sinal de \textit{wi-fi} esteja completamente preenchido e adote $\pi=3$. Para que as áreas das figuras hachuradas no setor circular e no retângulo sejam iguais, o percentual de bateria do celular será, aproximadamente, de:

\begin{enumerate}
	\item 81\%
	\item 86\%
	\item 91\%
	\item 96\%
\end{enumerate}
